% \subsection{Discussion}
This dissertation develops a framework for analyzing high-frequency repeated cross-sectional survey data that exhibit temporal correlation while containing individual-level covariates that are time independent. It does so by introducing a latent ARMA process, inducing an overall ARMA-Noise residual structure, and recursively fitting parameter estimates with GLS and ARMA-Noise methods. In order to accommodate survey data with varying daily sample size, this research furthered ARMA-Noise modeling by relaxing the preexisting condition that sample sizes be constant. Results provide insight into how temporal correlation interacts with cross-sectional noise, particularly how not properly accounting for this interaction can obfuscate results and hinder proper inference. Given how large RCS datasets can be, computational methods were tailored to the specific characteristics of this model by leveraging the relationship between the spectral decompositions of the various covariance matrices involved. The exploration of these models has produced insight into the limitations of their analysis and yielded several interesting directions for further research.
% \section{Interpretation of Key Findings}

A result central to this work is that parameter estimates for a latent ARMA process will be fundamentally flawed if fit to aggregate residuals, as the presence of additional white noise introduces bias in estimates for not only the variance parameter but the moving average parameter as well. Such an impact is tractable, however, and can be reverse engineered to obtain asymptotically unbiased estimates for the true time series parameters. The relation between the observable and latent time series must be incorporated into the analysis of time series estimate. This impact can change the order of the ARMA process on top of the parameters themselves. For instance, in the case that the latent time series is actually an AR(1) process, the observed aggregate errors will be ARMA(1,1) with a moving average parameter of opposite sign as that of the autoregressive parameter. Modeling the noisy time series as an AR(1) process would prevent a consistent estimate for $\phi$ from being obtained. Simulation results reinforce this conclusion. We see that as the relative variance of the additional noise increases, estimates for $\phi$ obtained by fitting an AR(1) model become increasingly biased, while those fit with an ARMA(1,1) model remain accurate.


% \section{Implications}
    This work provides a viable method for analyzing datasets such as the mVAM surveys, elucidating individual-level effects by properly accounting for temporal shocks that would otherwise obfuscate them. Such a method avoids the information loss inherent to the aggregate methods that are currently commonplace.

% \subsection{Methodological}


% \subsection{Practical}
Furthermore, failure to account for temporal correlation when modeling repeated cross-sectional data can lead to invalid inference. At the same time, flattening the data to be analyzed into aggregates results in significant information loss at the individual level. In particular, in the event that individual-level variables like education level or income source see little variance in their prevalence at the time-step level, fitting aggregated data using a time series with covariates could easily obscure an otherwise significant effect.

% \section{Limitations}
    This method is not without limitations, particularly given its reliance on asymptotic distributions. As such confidence regions underperform for low values of T, particularly in boundary cases like $\phi,\theta \approx \pm1$ and $\phi + \theta \approx 0$, as the resulting covariance structure is close to that of white noise. The imposition of an ARMA order no greater than $(1,1)$ is certainly limiting, as it precludes REMARCS Models with all but three types of latent ARMA processes from being properly fit.  

% \section{Directions for Future Research}
    There are a number of natural extensions that arise from this work, particularly 
    generalizing the time series level of the model to allow for $p,q > 1$. Given that 
    survey data so often tracks binary output variables, relaxing the Gaussian condition 
    is of further interest. Since confidence intervals are currently based on asymptotic
    distributions, the derivation of small-sample corrections is warranted. Lastly, 
    while the theory now accommodates samples size that vary, whether it be strictly stationary
    or not in certain circumstances, understanding the effect this has on parameter estimation
    is not yet understood.
% \section{Conclusion}